\documentclass[aspectratio=169]{beamer}

\usetheme{Madrid}
\usecolortheme{default}

\usepackage[utf8]{inputenc}
\usepackage{amsmath, amssymb}
\usepackage{graphicx}

\title{{name}}
\subtitle{A short subtitle}
\author{Your Name}
\institute{Your Institution}
\date{\today}

\begin{document}

\begin{frame}
    \titlepage
\end{frame}

\begin{frame}{Outline}
    \tableofcontents
\end{frame}

\section{Introduction}

\begin{frame}{The Problem}
    \begin{itemize}
        \item What the audience needs to care about.
        \item Why the obvious approach is not enough.
        \item What this talk adds.
    \end{itemize}
\end{frame}

\begin{frame}{Building Up a Point}
    % Items appear one at a time, so the audience reads with you
    % rather than ahead of you.
    \begin{itemize}
        \item<1-> First establish the setup.
        \item<2-> Then the complication.
        \item<3-> Then the resolution.
    \end{itemize}
\end{frame}

\section{Method}

\begin{frame}{Two Columns}
    \begin{columns}
        \begin{column}{0.5\textwidth}
            \begin{block}{Idea}
                A short statement of the approach.
            \end{block}
        \end{column}
        \begin{column}{0.5\textwidth}
            \begin{equation*}
                f(x) = \int_{0}^{x} g(t)\,dt
            \end{equation*}
        \end{column}
    \end{columns}
\end{frame}

\section{Results}

\begin{frame}{What Happened}
    \begin{enumerate}
        \item The headline result.
        \item The caveat that comes with it.
    \end{enumerate}
\end{frame}

\begin{frame}{Thank You}
    \centering
    \Large Questions?
\end{frame}

\end{document}
